\section{Expériences pour faire l'inventaire des résines}
Nous avons coulé des résines pour faire l'inventaire des propriété dont nous disposions. Nous en avons profité pour couler des bicouches pour étudier la qualité des interfaces. 

\subsection{Protocole}
\subsubsection{Protocole de coulée}


Liste des échantillons produits:
\subsubsection*{Comparaison des trois résines}
\begin{itemize}
	\item \textbf{ech1} — 1 bloc Epodex
	\item \textbf{ech2} — 1 bloc SR8200
	\item \textbf{ech3} — 1 bloc SR1670
\end{itemize}

\subsubsection*{Comparaison des proportions de durcisseur}
\begin{itemize}
	\item \textbf{ech4} — 1 bloc Epodex \(-30\)
	\item \textbf{ech5} — 1 bloc Epodex \(-15\)
	\item \textbf{ech6} — 1 bloc Epodex \(+15\)
	\item \textbf{ech7} — 1 bloc Epodex \(+30\)
\end{itemize}

\subsubsection*{Vérification interface}
\begin{itemize}
	\item \textbf{ech8} — 1 bicouche Epodex coulée en deux fois
	\item \textbf{ech9} — 1 bicouche Epodex / SR8200
	\item \textbf{ech10} — 1 bicouche Epodex / SR1670
\end{itemize}

\subsubsection*{Essai de charge}
\begin{itemize}
	\item \textbf{ech11} — 1 bloc SR8200 chargé en marbre
\end{itemize}
\vspace{1cm}
Détaillons les proportions des mélanges: 

\begin{itemize}
	\item \textbf{Epodex 0\%} mélange $R~=~197.5g \qquad D~=~99g$ avec nous avons remplie les échantillons e1,e8,e9 et e10. Les bicouches ont été remplis à moitier (53.9g), ils ont fuit nous avons sauvé l'expérience mais les nivaux sont à repérer. 
	\item \textbf{Epodex 30\%} mélange $R~=~67g \qquad D~=~51.5g$.
	\item \textbf{Epodex 15\%} mélange $R~=~73g \qquad D~=~45.5g$.
	\item \textbf{Epodex -30\%} mélange $R~=~91g \qquad D~=~28g$.
	\item \textbf{Epodex -15\%} mélange $R~=~85g \qquad D~=~33.5g$.
	\item \textbf{SR8200 + SD7403} mélange $R~=~92g \qquad D~=~34g$.
	\item \textbf{SR8200 + SD7403 chargé de marbre} mélange $R~=~92g \qquad D~=~34g$ ainsi que $30g$ de marbre.
	\item \textbf{SR1670 + SD7160} mélange $R~=~86.5g \qquad D~=~40.5g$.
\end{itemize}

Remarques: Il faudra mieux étanchéifier à l'avenir. Les prises rapides ont pas trop coulés. La charge à l'aire d'avoir marché. Les epodex ont des petites bulles.
\subsubsection{Protocole de mesures}

Il faut lister les mesures à faire.   


\subsection{Résultats}